% SEGMENT OLP-0698-S01
% Part: proof-theory
% Chapter: sequent-calculus
% Section: admissible-derivable

\documentclass[../../../include/open-logic-section]{subfiles}

\begin{document}

\olfileid{pt}{seq}{adm}

\olsection{அனுமதிக்கத்தக்க மற்றும் வருவிக்கத்தக்க விதிகள்}

ஒரு குறிப்பிட்ட விதியைக் கொண்டு !!{prove} செய்யக்கூடிய
அனைத்தையும் அந்த விதி இல்லாமலும் !!{prove} செய்ய
முடியுமா என்ற முடிவுகளை நிறுவல் கோட்பாட்டின் பல
முடிவுகள் பற்றியோ பயன்படுத்தியோ உள்ளன.
வெட்டு நீக்கத் தேற்றம் இதற்கு மிகவும் அறியப்பட்ட
எடுத்துக்காட்டு. \Cut{} உள்ள தொடரணி முறைமையின்
விதிகளைக் கொண்டு !!{prove} செய்யக்கூடிய எந்தத்
தொடரணியையும் \Cut இன்றியும் !!{prove} செய்யலாம்
என்பதை அது காட்டுகிறது. இத்தகைய பண்பு
முக்கியமானதால், அதற்கொரு பெயர் உண்டு.

\begin{defn}
ஒரு தொடரணி முறைமையுடன் விதி~$R$ ஐச் சேர்த்து
நிறுவக்கூடிய எந்தத் தொடரணியையும், $R$ இன்றி
அந்த முறைமையிலேயே நிறுவ முடிந்தால்,
அந்த முறைமையில் விதி~$R$ \emph{அனுமதிக்கத்தக்கது}.
\end{defn}

% SEGMENT OLP-0698-S02
\begin{ex}\ollabel{ex:land-deriv}
$\LeftR{\land}$ விதி \Log{G1c} யிலும்
\Log{G3c} யிலும் வேறுவேறு வடிவங்களைக் கொள்கிறது.
\Log{G1c} இல் அதற்கு இரு வடிவங்கள் உள்ளன:
ஒன்றின் முன்கூற்றில் முதன்மை
!!{formula}~$!A \land !B$ இன் இடக்கூறு~$!A$
உள்ளது; மற்றொன்றில் வலக்கூறு~$!B$ உள்ளது:
\[
\Axiom$ !A, \Gamma \fCenter \Delta$
\RightLabel{$\LeftR{\land}_1$}
\UnaryInf$ !A \land !B, \Gamma \fCenter \Delta$
\DisplayProof
\qquad
\Axiom$!B, \Gamma \fCenter \Delta$
\RightLabel{$\LeftR{\land}_2$}
\UnaryInf$!A \land !B, \Gamma \fCenter \Delta$
\DisplayProof
\]
இதற்கு மாறாக, \Log{G3c} இல் ஒரே ஒரு விதிதான் உண்டு:
\[
\Axiom$ !A, !B, \Gamma \fCenter \Delta$
\RightLabel{$\LeftR{\land}_3$}
\UnaryInf$ !A \land !B, \Gamma \fCenter \Delta$
\DisplayProof
\]
இப்போது “\Log{G3c} இன்
$\LeftR{\land}_3$ விதி \Log{G1c} இல்
அனுமதிக்கத்தக்கதா?” என்று கேட்கலாம்.
ஆம்: $\LeftR{\land}_3$ ஐப் பயன்படுத்தும்
எந்த அனுமானத்தையும் \Log{G1c} இன்
விதிகளால் நேரடியாகப் பாவனை செய்யலாம்.
$\LeftR{\land}_1$, $\LeftR{\land}_2$
ஆகியவற்றோடு, கட்டமைப்பு விதியான
$\LeftR{\Contraction}$ உம் தேவைப்படும்:
\[
\Axiom$!A, !B, \Gamma \fCenter \Delta$
\RightLabel{$\LeftR{\land}_1$}
\UnaryInf$!A \land !B, !B, \Gamma \fCenter \Delta$
\RightLabel{$\LeftR{\land}_2$}
\UnaryInf$!A \land !B, !A \land !B, \Gamma \fCenter \Delta$
\RightLabel{\LeftR{\Contraction}}
\UnaryInf$!A \land !B, \Gamma \fCenter \Delta$
\DisplayProof
\]
\end{ex}

% SEGMENT OLP-0698-S03
ஒரு விதியைப் பயன்படுத்தும் அனுமானங்களை
மற்ற விதிகளால் பாவனை செய்வது, அந்த விதி
அனுமதிக்கத்தக்கது என்பதைக் காட்டும் ஒரு வழி.
ஆனால் அதுவே ஒரே வழியல்ல; இவ்வாறு பாவனை
செய்ய முடியாவிட்டாலும் சில விதிகள்
அனுமதிக்கத்தக்கவை. பாவனை செய்ய முடியும்
விதிகளை \emph{வருவிக்கத்தக்கவை} என்போம்.

\begin{defn}
ஒரு விதி~$R$,
\[
\AxiomC{$S_1$}
\AxiomC{$\dots$}
\AxiomC{$S_n$}
\RightLabel{$R$}
\TrinaryInfC{$S$}
\DisplayProof
\]
என்பது, முறைமையின் விதிகளையும்
அடிக்கோள்களையும், மேலும் $S_1$, \dots,~$S_n$
வடிவங்களிலான கூடுதல் தொடக்கத் தொடரணிகளையும்
கொண்டு $S$ தொடரணிக்கு ஒரு திட்டவடிவ
!!{proof} இருந்தால், அந்த முறைமையில்
\emph{வருவிக்கத்தக்கது} எனப்படும்.
\end{defn}

\olref{ex:land-deriv} இல் தரப்பட்ட !!{proof},
$\LeftR{\land}_3$ என்பது \Log{G1c} இல்
வருவிக்கத்தக்க விதி என்பதைக் காட்டுகிறது.
ஏனெனில் அந்த !!{proof},
$\LeftR{\land}_3$ இன்
முன்கூற்றுகளைத் தொடக்கத் தொடரணிகளாகக் கொண்டு,
\Log{G1c} இன் விதிகளை மட்டும் பயன்படுத்தி
$\LeftR{\land}_3$ இன் முடிவுக்கான திட்டவடிவ
!!{proof} ஆகிறது. (\Log{G1c} இன்
அடிக்கோள்கள் எதையும் அது பயன்படுத்தவில்லை.)

\begin{prop}\ollabel{prop:lor-G3-adm}
\Log{G3c} இன் $\LeftR{\land}$,
$\RightR{\lor}$ விதிகள் \Log{G1c} இல்
வருவிக்கத்தக்கவை.
\end{prop}

\begin{proof}
  $\LeftR{\land}$ க்கான நிறுவல்
  \olref{ex:land-deriv} இல் உள்ளது.
  $\RightR{\lor}$ க்கான நிறுவல்
  பயிற்சியாக விடப்படுகிறது.
\end{proof}

\begin{prob}
\Log{G3c} இன் $\RightR{\lor}$ விதி
\Log{G1c} இல் வருவிக்கத்தக்கது
என்பதைக் காட்டுக.
\end{prob}

\begin{prob}
$\liff$ என்பது வரையறுக்கப்பட்ட
!!{operator} என்று கொள்க:
$!A \liff !B$ என்பது
$(!A \lif !B) \land (!B \lif !A)$
என்பதன் சுருக்கக் குறியீடு.
பின்வரும் விதிகள் (a)~\Log{G1c} யிலும்
(b)~\Log{G3c} யிலும் வருவிக்கத்தக்கவை
என்பதைக் காட்டுக:
\begin{enumerate}
\item \Axiom$\Gamma, !A, !B \fCenter \Delta$
\Axiom$\Gamma \fCenter \Delta, !A, !B$
\RightLabel{\LeftR{\liff}}
\BinaryInf$!A \liff !B, \Gamma \fCenter \Delta$
\DisplayProof
\item
\Axiom$!A, \Gamma \fCenter \Delta, !B$
\Axiom$!B, \Gamma \fCenter \Delta, !A$
\RightLabel{\RightR{\liff}}
\BinaryInf$\Gamma \fCenter \Delta, !A \liff !B$
\DisplayProof
\end{enumerate}
\end{prob}

% SEGMENT OLP-0698-S04
\begin{ex}\ollabel{ex:weak-adm}
\Log{G1c} இன் $\RightR{\Weakening}$ விதி,
\[
\Axiom$ \Gamma \fCenter \Delta$
\RightLabel{\RightR{\Weakening}}
\UnaryInf$ \Gamma \fCenter \Delta, !A$
\DisplayProof
\]
\Log{G3c} இல் அனுமதிக்கத்தக்கது.
நிறுவலின் கருத்து எளிது: \Log{G3c} இல்
முன்கூற்று $\Gamma \Sequent \Delta$
வுக்கான !!a{proof}~$\pi$ உள்ளது
என்று கொள்வோம். $\pi$ இன் ஒவ்வொரு
தொடரணியின் வலப்புறத்திலும்
!!{formula}~$!A$ ஐச் சேர்ப்போம்.
அடிக்கோள்கள் அடிக்கோள்களாகவே இருக்கும்;
சரியான அனுமானங்களும் சரியானவையாகவே
இருக்கும். புதிய !!{proof} இன்
இறுதித் தொடரணி $\Gamma \Sequent
\Delta, !A$ ஆகும்.

% SEGMENT OLP-0698-S05
ஆனால் $\RightR{\Weakening}$ விதி
\Log{G3c} இல் வருவிக்கத்தக்கதல்ல.
அது வருவிக்கத்தக்கது என்றால்,
\Log{G3c} இன் அடிக்கோள்களையும்
விதிகளையும், தேவையானால் கூடுதல்
தொடக்கத் தொடரணியான
$\Gamma \Sequent \Delta$ ஐயும்
பயன்படுத்தி $\Gamma \Sequent
\Delta, !A$ க்கான !!a{proof} ஐக்
கண்டுபிடிக்கலாம் என்று கொள்வோம்.
$\RightR{\Weakening}$ வருவிக்கத்தக்கதாக
இருக்க வேண்டியதுபோல் அந்த !!{proof}
முழுமையான திட்டவடிவில் இருந்தால்,
$\Gamma$ வையும்~$\Delta$ வையும்
நீக்கி அதை $\Sequent !A$ என்பதற்கான
!!a{proof} ஆக மாற்றலாம். அப்போது
கூடுதல் தொடக்கத் தொடரணிகளான
$\Gamma \Sequent \Delta$,
வெற்றுத் தொடரணி~$\quad \Sequent \quad$
ஆக மாறும். வெற்றுத் தொடரணியில்
!!{formula}s எதுவுமே இல்லை;
ஆனால் \Log{G3c} இன் எல்லா விதிகளின்
முன்கூற்றுகளிலும் குறைந்தது ஒரு செயலில்
உள்ள !!{formula} தேவைப்படுவதால்,
இந்தத் திட்டவடிவ !!{proof} இல்
வெற்றுத் தொடரணி எந்த அனுமானத்துக்கும்
முன்கூற்றாக முடியாது. ஆகவே கூடுதல்
தொடரணி $\Gamma \Sequent \Delta$ ஐ
உண்மையில் தொடக்கத் தொடரணியாகப்
பயன்படுத்தாதிருந்தால்தான் அந்த
!!{proof} திட்டவடிவமாக இருக்க முடியும்.
வேறு சொற்களில், அது \Log{G3c} இன்
அடிக்கோள்களையும் விதிகளையும் மட்டும்
பயன்படுத்தி $\Gamma \Sequent
\Delta, !A$ க்கான திட்டவடிவ
வருவித்தலாக இருக்கும்; அந்த
வடிவிலான எந்தத் தொடரணிக்கும்
\Log{G3c} இல் !!a{proof}
இருக்கிறது என்று காட்டியிருப்போம்.
ஆனால் அது இயலாது: \Log{G3c}
சரியானது; ஆகவே ஒவ்வொரு
!!{structure} இலும் மெய்யான
தொடரணிகளையே அது !!{prove} செய்யும்.
$\Gamma = \Delta
= \emptyset$ என்றும்
$!A \ident \lfalse$ என்றும் எடுத்தால்,
எடுத்துக்காட்டாக, $\quad \Sequent \lfalse$
என்ற தொடரணியை \Log{G3c}
!!{prove} செய்ய வேண்டியிருக்கும்;
அது அவ்வாறு செய்யாது.
\end{ex}

% SEGMENT OLP-0698-S06
\Log{G3c} இல் தளர்த்தல் அனுமதிக்கத்தக்க விதி
என்பதற்கு ஓர் அழகிய தொகுத்தறிதல் நிறுவலைத்
தருவோம். உண்மையில், உள்ளுணர்வு வாதம்
காட்டுவதுபோல், சற்றுக் கூடுதலான வலிமையுள்ள
முடிவு கிடைக்கிறது: $\pi$ இன் ஒவ்வொரு
தொடரணியின் வலப்புறத்திலும் $!A$ ஐச்
சேர்க்கும்போது, !!{proof} இன் கட்டமைப்பு
மாறுவதில்லை; குறிப்பாக, அதன் உயரம்
அதிகரிப்பதில்லை. இது முக்கியமானதால்,
இதற்குத் தனி வரையறை தருவோம்.

\begin{defn}
ஒரு விதி~$R$,
\[
\AxiomC{$S_1$}
\AxiomC{$\dots$}
\AxiomC{$S_n$}
\RightLabel{$R$}
\TrinaryInfC{$S$}
\DisplayProof
\]
என்பது, $i = 1$, \dots,~$n$ ஒவ்வொன்றுக்கும்
$\Proves[h_i] S_i$ எனில்,
$m = \max(h_1, \dots,
h_n)$ உடன் $\Proves_m S$
ஆகுமானால், அந்த முறைமையில்
\emph{!!{height}-பேணும் அனுமதிக்கத்தக்க}
(அல்லது \emph{!!{hp}-அனுமதிக்கத்தக்க})
விதி எனப்படும்.
\end{defn}

ஒரு விதி !!{height}-பேணும்
அனுமதிக்கத்தக்க விதி என்பதை நிறுவ,
முன்கூற்றுகள்~$S_i$ க்கு
!!{height}~$h_i = \pheight{\pi_i}$
உடைய !!{proof}s~$\pi_i$ இருக்கும்போதெல்லாம்,
முடிவுக்கு $\pi$ என்ற !!a{proof}
ஒன்று இருந்து, அதன்
!!{height}~$\pheight{\pi}$,
$h_i$ களின் மிகப்பெரிய மதிப்பை
மீறாது என்பதைக் காட்டினால் போதும்.
குறிப்பாக, ஒரே ஒரு முன்கூற்று~$S_1$
மட்டுமே இருந்தால்,
$\pheight{\pi} \le \pheight{\pi_1}$
என்பதே போதும்.
$\RightR{\Weakening}$ மற்றும்
$\LeftR{\Weakening}$ நிலைகளில்
உண்மையில் $\pheight{\pi} = \pheight{\pi_1}$
கிடைக்கிறது; ஆனால் அது தேவையானதல்ல.

% SEGMENT OLP-0698-S07
\begin{prop}\ollabel{prop:weak-G3c-adm}
\Log{G1c} இன் $\RightR{\Weakening}$,
$\LeftR{\Weakening}$ விதிகள் \Log{G3c} இல்
!!{height}-பேணும் அனுமதிக்கத்தக்க விதிகள்.
\end{prop}

\begin{proof}
  $\pheight{\pi} = n$ என்ற
  !!{height} உடைய !!{proof}~$\pi$
  கொண்டு $\Log{G3c} \Proves \Gamma \Sequent \Delta$
  எனில், $\pheight{\pi'} = n$ உடைய
  $\Gamma \Sequent \Delta, !A$ என்பதன்
  \Log{G3c}-!!{proof}~$\pi'$ உண்டு
  என்பதைக் காட்ட வேண்டும். $n$ மீது
  தொகுத்தறிந்து நிறுவுவோம். $n=0$
  என்றால், $\pi$ என்பது
  $\Gamma \Sequent \Delta$ என்ற
  அடிக்கோள்தான்; அதாவது ஓர் அணு
  $!C$, $\Gamma$ விலும்~$\Delta$
  விலும் வருகிறது. ஆகவே
  $\Gamma \Sequent \Delta, !A$ உம்
  அடிக்கோள்தான்.

% SEGMENT OLP-0698-S08
  $n = \pheight{\pi} >0$ என்றால்,
  !!{proof}~$\pi$ ஒரு கடைசி
  அனுமானத்தைக் கொண்டிருக்கும்.
  எந்த விதி பயன்படுத்தப்பட்டது
  என்பதைப் பொறுத்து நிலைகளைப்
  பிரிக்கிறோம். ஓர் எடுத்துக்காட்டை
  மட்டும் தருகிறோம்: கடைசி
  அனுமானம் $\RightR{\land}$
  என்று கொள்வோம். அப்போது
  $\Delta = \Delta', !B \land !C$;
  கடைசி அனுமானத்தின் உடனடித்
  துணை-!!{proof}s $\pi_1$ உம்~$\pi_2$
  உம் முறையே $\Gamma \Sequent \Delta', !B$,
  $\Gamma \Sequent
  \Delta', !C$ என்ற இறுதித்
  தொடரணிகளைக் கொண்டிருக்கும்.
  அதாவது, $\pi$ இவ்வாறு இருக்கும்:
  \[
  \AxiomC{}
  \RightLabel{$\pi_1$}
  \Deduce$\Gamma \fCenter \Delta', !B$
  \AxiomC{}
  \RightLabel{$\pi_2$}
  \Deduce$\Gamma \fCenter \Delta', !C$
  \RightLabel{\RightR{\land}}
  \BinaryInf$\Gamma \fCenter \Delta', !B \land !C$
  \DisplayProof
  \]
  மேலும் $n =
  \max(\pheight{\pi_1}, \pheight{\pi_2}) + 1$.
  எனவே $\pheight{\pi_1} <
  n$ மற்றும் $\pheight{\pi_2} < n$;
  தொகுத்தறிதல் கருதுகோள் பொருந்தும்:
  $\Gamma \Sequent
  \Delta', !B, !A$ மற்றும்
  $\Gamma \Sequent \Delta', !C, !A$
  என்பவற்றுக்கான !!{proof}s
  $\pi_1'$ உம்~$\pi_2'$ உம் கிடைக்கும்.
  $\RightR{\land}$ விதியால்,
  $\Gamma \Sequent
  \Delta', !B \land !C, !A$
  என்பதன் !!a{proof}~$\pi'$
  ஐப் பெறுகிறோம்:
  \[
  \AxiomC{}
  \RightLabel{$\pi_1'$}
  \Deduce$\Gamma \fCenter \Delta', !B, !A$
  \AxiomC{}
  \RightLabel{$\pi_2'$}
  \Deduce$\Gamma \fCenter \Delta', !C, !A$
  \RightLabel{\RightR{\land}}
  \BinaryInf$\Gamma \fCenter \Delta', !B \land !C, !A$
  \DisplayProof
  \]
  இதன் உயரம்
  $\pheight{\pi'} =
  \max(\pheight{\pi_1'}, \pheight{\pi_2'}) + 1 = \max(\pheight{\pi_1},
  \pheight{\pi_2}) + 1 = \pheight{\pi}$.
\end{proof}

% SEGMENT OLP-0698-S09
\olref{prop:weak-G3c-adm} மிகவும் பயனுள்ளது.
இதைத் திரும்பத் திரும்பப் பயன்படுத்துவதற்கு
ஒரு குறியீட்டை அறிமுகப்படுத்துவோம்:
$\pi$ என்பது~$\Gamma \Sequent \Delta$
என்பதன் !!a{proof} என்றால்,
$\pi[\Pi \Sequent \Lambda]$
என்பது, இடப்புறத்தில்~$\Pi$ இல் உள்ள
எல்லா !!{formula}s உடனும்,
வலப்புறத்தில்~$\Lambda$ இல் உள்ள
எல்லா !!{formula}s உடனும்
தளர்த்தலின் அனுமதிக்கத்தன்மையைப்
பயன்படுத்தியதன் விளைவு.
இது $\Gamma, \Pi
\Sequent \Delta, \Lambda$
என்பதன் !!a{proof} ஆகும்.

\end{document}
